% %%%%%%%%%%%%%%%%%%%%%%%%%%%%%%%%%%%%%%%%%%%%%%%%%%%%%%%%%%%%%%%%%%%%%
% %% Theory %%
% GRC
% %%%%%%%%%%%%%%%%%%%%%%%%%%%%%%%%%%%%%%%%%%%%%%%%%%%%%%%%%%%%%%%%%%%%%

\section{Isolating the Input Reads in a Generic Cache}
\label{sec:theory:isolate_input_reads}

\lettrine{W}{hen} computing an \acrshort{op} \acrshort{spmv}, all the \glspl{irrAccess} come from updating the output vector $c$.
However, the input matrix $A$ and vector $b$ are read sequentially, thus being efficiently cached by a generic cache that naturally exploits spatial locality.

With the matrix $A$ stored in \acrshort{csc} format, each iteration of the \acrshort{op} algorithm necessitates reading four tables: $\mathit{A.ptr}$, $\mathit{A.idx}$, $\mathit{A.val}$ and $b$.
In the same iteration, a fifth table is updated ``randomly'': $c$.
If we consider cache-lines of size $l$ elements, then each cache-line read of the $\mathit{A.ptr}$ table brings $l$ consecutive entries, which are pointers to $l$ consecutive columns of the matrix.
Similarly, each cache-line read of the $b$ table brings $l$ consecutive entries, which will be multiplied by the non-zero elements of $l$ consecutive columns of the matrix.
These reads are needed in the outer-loop of the \acrshort{op} algorithm.
In the inner-loop, each cache-line read of the $\mathit{A.idx}$ and $\mathit{A.val}$ tables brings $l$ consecutive non-zero elements of the matrix to process, which will generate $l$ ``random'' updates to the output vector $c$, thus $l$ entire cache-lines of $c$ may be accessed in the worst case.

To avoid unnecessary \glspl{eviction}, at each iteration, the generic cache should obviously keep $\mathit{A.ptr}$, $\mathit{A.idx}$, $\mathit{A.val}$ and $b$ cache-lines up until they are no longer needed.
That means $l$ iterations for $\mathit{A.idx}$ and $\mathit{A.val}$, and around $l \times \mathit{nnz}_c$ iterations for $\mathit{A.ptr}$ and $b$.
However, the cache must also keep the cache-lines of $c$ that are being updated, which, over the course of $l \times \mathit{nnz}_c$ iterations, can necessitate up to the same amount of different cache-lines in the worst case, on unpredictable addresses.
Thus, it would be difficult for a generic cache to keep all the useful data until it is no longer needed.
Especially $\mathit{A.ptr}$ and $b$ cache-lines, which are accessed once every $\mathit{nnz}_c$ iterations, making \gls{eviction} policies such as \acrshort{lru} unaware of their future use.

By isolating the input reads in a dedicated cache, as proposed with \acrlong{spdc} and the \gls{redC} baseline in the previous chapter, we greatly increase the chances that all the cache-lines of $A$ and $b$ remain in the cache until they are no longer needed.

With a dedicated cache region for the inputs, the \acrshort{grc} in previous chapter, only the four sequentially read tables need to be considered.
Each of these tables need to be read only once.
With a classical \acrshort{lru} policy, $\mathit{A.ptr}$ and $b$ cache-lines have lower chances of being evicted since the only accesses occurring while processing a column are the $\sim \frac{\mathit{nnz}_c}{l}$ cache-line reads of $\mathit{A.idx}$ and $\mathit{A.val}$.
$\mathit{A.idx}$ and $\mathit{A.val}$ themselves can not be undesirably evicted anymore.

To avoid evicting $\mathit{A.ptr}$ and $b$ too soon, the \acrshort{grc} must have enough size to handle the $\frac{\mathit{nnz}_c}{l}$ cache-line reads while computing a column.
Thus, we define in Equation~\ref{eq:theory:grc_size} the minimum size of the \acrshort{grc} in elements, which would be needed to achieve best performance in memory traffic.

\begin{equation}
    s_{C_\mathit{min}}(\mathit{GRC}) = 2 \times l + 2 \times \mathit{nnz}_c
    \label{eq:theory:grc_size}
\end{equation}

This relationship is not exact, as it does not consider address alignment effects, but it gives a good approximation of the minimum size needed for the \acrshort{grc} to avoid unnecessary \glspl{eviction}.
It explains why, in the previous chapter, the total memory traffic slightly drops when the \acrshort{grc} size is decreased.

Separating the inputs in a dedicated cache region also has the advantage of allowing prefetching techniques to be used efficiently.
It has no impact on memory traffic but can improve performance by hiding memory latencies.

\begin{highlight}{Conclusion}{Red}
    To avoid unnecessary \glspl{eviction}, it is profitable to \textbf{separate} the sequential reads of the inputs $A$ and $b$ from the ``random'' updates of the output $c$, into \textbf{dedicated cache regions}.
    A generic cache of reasonable size can achieve \textbf{low memory traffic} when processing $A$ and $b$.
\end{highlight}